% !TEX program  = xelatex
% !TEX root     = coursework_imu.tex
% !TEX root = coursework_imu.tex
\documentclass[a4paper, 14pt]{extarticle}

\usepackage{fontspec}

% Основной шрифт — Times New Roman
\setmainfont{Times New Roman}
\setmonofont{Liberation Mono}


\usepackage{listings}
\lstset{
    basicstyle = \ttfamily,
    numbers = left,
    numbersep = 5pt,
    tabsize = 4,
    showstringspaces = false,
    keepspaces = true
}
\usepackage{amsmath}
\usepackage{enumitem}
\usepackage{indentfirst}
%\usepackage[T2A]{fontenc}
\usepackage[english, russian]{babel}
\usepackage{multirow}
\usepackage{tabularx}

\usepackage{subcaption}

\usepackage{array}
\usepackage{booktabs}
\usepackage{colortbl}
% Центрируем содержимое multirow по вертикали
\usepackage[table]{xcolor}
\usepackage{tikz}
\usetikzlibrary{backgrounds}

% удобный тип колонки с фиксированной шириной и центрированием
\newcolumntype{P}[1]{>{\centering\arraybackslash}p{#1}}

% малый корректор (если нужно подровнять край — поставьте 0.4pt или 0.2pt)
\newlength{\diagfudge}
\setlength{\diagfudge}{0pt}

% #1 = цвет нижне-левого треугольника
% #2 = цвет верхне-правого треугольника
% #3 = содержимое ячейки (одно значение)
\newcommand{\diagcell}[3]{%
  \begingroup
  \setlength{\fboxsep}{0pt}% убрать лишние отступы вокруг коробок
  % рисуем фон в "нулевой" коробке, сдвинутой влево на \tabcolsep
  \makebox[0pt][l]{\hspace*{-\tabcolsep}%
    \tikz[baseline=(box.base)]{%
      \node[inner sep=0pt, outer sep=0pt, anchor=west,
            minimum width=\dimexpr\linewidth+2\tabcolsep+\diagfudge\relax,
            minimum height=1.5\baselineskip] (box) {};
      \begin{scope}[on background layer]
        \fill[#1]  (box.south west) -- (box.south east) -- (box.north west) -- cycle;
        \fill[#2]  (box.north east) -- (box.south east) -- (box.north west) -- cycle;
        % Для отладки раскомментируйте следующую строку, чтобы увидеть рамку узла:
        % \draw[red,thin] (box.north west) rectangle (box.south east);
      \end{scope}
    }%
  }%
  % Теперь выводим само содержимое, ровно в ширине \linewidth и по центру
  \makebox[\linewidth][c]{\strut #3}%
  \endgroup
}


\usepackage{textcomp}
\newcommand{\degree}{\ensuremath{{}^\circ}} %использвоание значка градуса в $$

\usepackage{geometry}
\geometry{top=20mm}
\geometry{bottom=20mm}
\geometry{left=30mm}
\geometry{right=15mm}


\usepackage{titlesec}
\titleformat{\section}{\normalsize\it\centering}{\thesection}{1em}{}
\titleformat{\subsection}{\normalsize\it}{\thesubsection}{1em}{}

%\titleformat{\subsection}
%  {\normalfont\normalsize\bfseries\centering}
%  {\thesubsection}{1em}{}


\usepackage{tocloft}

%\renewcommand{\cftsecaftersnum}{.} % Добавляет точки после номера

% Заменяем лидеры на точки для всех уровней
\renewcommand{\cftdotsep}{1}
\renewcommand{\cftsecleader}{\cftdotfill{\cftdotsep}}
\renewcommand{\cftsubsecleader}{\cftdotfill{\cftdotsep}}
\renewcommand{\cftsubsubsecleader}{\cftdotfill{\cftdotsep}}

% Ваши текущие настройки
\tolerance=1000
\emergencystretch=3em
\pretolerance=150
\usepackage[none]{hyphenat}

% Добавьте эти строки для подавления предупреждений
\hbadness=10000  % не показывать предупреждения об underfull
\vbadness=10000  % то же для вертикального направления


\usepackage{graphicx}
\graphicspath{ {./figures/} }

\usepackage{caption}
\captionsetup[figure]{labelsep=endash}

\usepackage[
    %hidelinks  % ссылки без цвета и рамок, но остаются кликабельными
]{hyperref}


\usepackage{newtxmath}

\usepackage{caption}
\captionsetup[table]{justification=raggedright,singlelinecheck=false,labelsep=endash}

\usepackage{float}


 
\usepackage{amssymb}
\usepackage{bm}
\usepackage{mathtools}
\usepackage{cite}
\usepackage{setspace}
 
\lstset{
    basicstyle        = \ttfamily\small,
    numbers           = left,
    numbersep         = 5pt,
    tabsize           = 4,
    showstringspaces  = false,
    keepspaces        = true,
    breaklines        = true,
    frame             = single,
    captionpos        = b,
    language          = C++
}
 
\setlength{\tabcolsep}{6pt}
 
% ─────────────────────────────────────────────────────────────────────────────
\begin{document}
\linespread{1.5}\selectfont
 
% ══════════════════════════ ТИТУЛЬНЫЙ ЛИСТ ══════════════════════════════════
\thispagestyle{empty}
\begin{center}
  Санкт-Петербургский политехнический университет Петра Великого\\[2pt]
  Институт компьютерных наук и технологий\\[2pt]
  Высшая школа программной инженерии
\end{center}
 
\vfill
 
\begin{center}
  \textbf{Курсовая работа}\\[4pt]
  \textbf{по дисциплине <<Объектно-ориентированное программирование>>}
\end{center}
 
\vskip8mm
Тема: <<Система визуализации инерциальной навигации на основе
полётного контроллера Pixhawk>>
\vskip10mm
 
Выполнил\\
Студент гр.~3331506/30401 \hfill Томилов~А.\,С.
\vskip8mm
Преподаватель \hfill Ананьевский~М.\,С.
\vskip10mm
\hfill <<\_\_\_>>\_\_\_\_\_\_\_\_\_ 2026~г.
 
\vfill
\begin{center}
  Санкт-Петербург\\
  2026~г.
\end{center}
\newpage
 
% ══════════════════════════ СОДЕРЖАНИЕ ══════════════════════════════════════
\tableofcontents
\newpage
 
% ─────────────────────────────────────────────────────────────────────────────
\section*{Введение}
\addcontentsline{toc}{section}{Введение}
% ─────────────────────────────────────────────────────────────────────────────
 
Инерциальная навигация~--- один из фундаментальных методов определения
положения и ориентации объекта в пространстве без использования внешних
радионавигационных сигналов~\cite{woodman2007}. В основе метода лежат данные
инерциального измерительного блока (IMU~--- \textit{Inertial Measurement Unit}),
включающего трёхосевые акселерометр и гироскоп. Интегрируя угловые скорости и
линейные ускорения во времени, можно восстановить полную навигационную картину:
ориентацию, скорость и положение объекта~\cite{titterton2004}.
 
Полётный контроллер Pixhawk является широко распространённым стандартом в области
автономных летательных аппаратов~\cite{px4_docs}. Он оснащён высококачественным
МЭМС-IMU и поддерживает протокол MAVLink~\cite{mavlink_proto}~--- открытый
телеметрический стандарт, позволяющий в реальном времени получать сырые данные
датчиков на внешний компьютер~\cite{mavlink_docs}.
 
Целью данной работы является разработка программной системы, которая осуществляет
чтение сырых IMU-данных (сообщения \texttt{RAW\_IMU} протокола MAVLink) с
полётного контроллера Pixhawk и выполняет их обработку в реальном времени:
вычисление ориентации объекта через кватернионное интегрирование
гироскопа~\cite{kuipers1999}, а также оценку скорости и позиции посредством
двойного интегрирования ускорения~\cite{groves2013}. Результаты визуализируются
в интерактивном 3D-окне с использованием OpenGL~3.3 и библиотеки
GLFW~\cite{glfw_docs}.
 
\textbf{Задачи работы:}
\begin{enumerate}[leftmargin=2cm]
    \item реализовать автономный модуль первичного считывания данных по протоколу
          MAVLink через последовательный порт с записью в CSV;
    \item разработать многопоточную архитектуру с потокобезопасным разделяемым
          буфером;
    \item реализовать алгоритм калибровки и кватернионного интегрирования
          гироскопических данных;
    \item реализовать алгоритм двойного интегрирования акселерометра для
          оценки позиции;
    \item реализовать механизм ZUPT (\textit{Zero Velocity Update}) для
          частичной компенсации дрейфа;
    \item создать OpenGL-визуализацию с 3D-моделью полётного контроллера,
          траекторией движения и HUD-панелью графиков датчиков;
    \item провести анализ точности полученных результатов.
\end{enumerate}
 
\newpage
 
% ─────────────────────────────────────────────────────────────────────────────
\section{Первичное считывание данных с Pixhawk}
% ─────────────────────────────────────────────────────────────────────────────
 
\subsection{Протокол MAVLink и сообщение RAW\_IMU}
 
MAVLink (\textit{Micro Air Vehicle Link})~--- лёгкий бинарный протокол
сериализации для беспилотных систем, разработанный Лоренцем Майером
в 2009~году~\cite{mavlink_docs}. Протокол широко используется в экосистемах
PX4 и ArduPilot~\cite{px4_docs}. Для получения сырых показаний датчиков
применяется сообщение \texttt{RAW\_IMU} (идентификатор~27 в диалекте
\texttt{common}), содержащее девять целочисленных полей в единицах внутренней
шкалы датчика (LSB~--- \textit{Least Significant Bit}):
 
\begin{table}[H]
\centering
\caption{Поля сообщения MAVLink \texttt{RAW\_IMU}}
\label{tab:rawimu}
\begin{tabularx}{\textwidth}{|>{\centering\arraybackslash}p{4.5cm}|
                               >{\centering\arraybackslash}p{3cm}|
                               >{\centering\arraybackslash}X|}
\hline
\textbf{Поле} & \textbf{Датчик} & \textbf{Описание} \\
\hline
\texttt{xacc, yacc, zacc}   & Акселерометр & Ускорение по осям X, Y, Z (LSB) \\
\hline
\texttt{xgyro, ygyro, zgyro}& Гироскоп     & Угловая скорость по осям (LSB) \\
\hline
\texttt{xmag, ymag, zmag}   & Магнитометр  & Магнитное поле по осям (LSB) \\
\hline
\end{tabularx}
\end{table}
 
\subsection{Модуль первичного считывания}
 
На первом этапе работы был написан автономный модуль
\texttt{read\_raw\_data\_from\_pixhawk.cpp}, позволивший проверить
корректность соединения с Pixhawk и собрать сырые данные для анализа и
подбора калибровочных коэффициентов. Взаимодействие реализовано через
последовательный порт \texttt{/dev/ttyACM0} на скорости 115\,200~бод.
Инициализация порта выполнена через POSIX-интерфейс
\texttt{termios}~\cite{posix_termios}: 8~бит данных, без паритета, один
стоп-бит, без аппаратного управления потоком, таймаут~1~с.
 
\begin{lstlisting}[caption={Настройка последовательного порта через termios}]
int setup_serial_port(const char* port, int baudrate) {
    int fd = open(port, O_RDWR | O_NOCTTY);
    struct termios tty;
    memset(&tty, 0, sizeof(tty));
    tcgetattr(fd, &tty);
    cfsetospeed(&tty, baudrate);
    cfsetispeed(&tty, baudrate);
    tty.c_cflag |= (CLOCAL | CREAD);
    tty.c_cflag &= ~CSIZE;
    tty.c_cflag |= CS8;
    tty.c_cflag &= ~(PARENB | CSTOPB | CRTSCTS);
    tty.c_lflag &= ~(ICANON | ECHO | ECHOE | ISIG);
    tty.c_iflag &= ~(IXON | IXOFF | IXANY);
    tty.c_oflag &= ~OPOST;
    tty.c_cc[VMIN]  = 0;
    tty.c_cc[VTIME] = 10;
    tcsetattr(fd, TCSANOW, &tty);
    return fd;
}
\end{lstlisting}
 
MAVLink-парсер реализован в виде конечного автомата:
\texttt{mavlink\_parse\_char()} возвращает ненулевое значение только
при получении полного валидного пакета с проверкой CRC-16~\cite{mavlink_proto}.
Каждый принятый пакет \texttt{RAW\_IMU} записывается в CSV-файл с меткой
времени с точностью до микросекунды~(\texttt{gettimeofday()}); принудительный
сброс буфера выполняется каждые 50~пакетов.
 
\subsection{Результаты первичного считывания}
 
Сообщения \texttt{RAW\_IMU} стабильно поступали с частотой порядка
100~Гц, что соответствует штатной конфигурации IMU на
Pixhawk~\cite{px4_docs}. Сырые данные акселерометра в покое показали
устойчивое значение по оси~Z (проекция силы тяжести), а гироскоп~---
дрейф нуля, подлежащий компенсации на следующем этапе.
 
\newpage
 
% ─────────────────────────────────────────────────────────────────────────────
\section{Архитектура системы визуализации}
% ─────────────────────────────────────────────────────────────────────────────
 
\subsection{Общая структура}
 
На основе опыта, полученного при реализации модуля чтения, была
разработана полноценная система \textbf{Pixhawk IMU Visualizer}~---
многопоточное приложение на C++14. Два параллельных потока взаимодействуют
через потокобезопасный разделяемый буфер (\texttt{SharedBuffer}):
 
\textit{Главный поток} (\texttt{pixhawk\_imu\_reader.cpp}) устанавливает
соединение с Pixhawk по UART, непрерывно парсит входящий байтовый поток
и при получении пакета \texttt{RAW\_IMU} помещает данные в очередь, а также
дублирует их в CSV-файл.
 
\textit{Поток визуализации} (\texttt{imu\_visualizer.cpp}) извлекает пакеты
из очереди, запускает алгоритм инерциальной навигации и обновляет 3D-сцену на
каждом кадре.
 
\subsection{Структура файлов проекта}
 
\begin{table}[H]
\centering
\caption{Файлы проекта}
\label{tab:files}
\begin{tabularx}{\textwidth}{|l|X|}
\hline
\textbf{Файл} & \textbf{Назначение} \\
\hline
\texttt{imu\_shared.h}            & Общие структуры данных: \texttt{RawIMU},
                                    \texttt{NavState}, \texttt{SharedBuffer},
                                    константы калибровки \\
\hline
\texttt{imu\_visualizer.h}        & Объявление функции \texttt{run\_visualizer()} \\
\hline
\texttt{imu\_visualizer.cpp}      & Навигационный алгоритм, GLSL-шейдеры,
                                    геометрия модели и OpenGL-рендер \\
\hline
\texttt{pixhawk\_imu\_reader.cpp} & MAVLink-ридер, запуск потока визуализации \\
\hline
\texttt{read\_raw\_data\_from\_pixhawk.cpp} & Автономный модуль первичного считывания \\
\hline
\texttt{CMakeLists.txt}           & Система сборки CMake \\
\hline
\texttt{c\_library\_v2/}          & Заголовочная библиотека MAVLink~C~\cite{mavlink_docs} \\
\hline
\end{tabularx}
\end{table}
 
\subsection{Разделяемые структуры данных (\texttt{imu\_shared.h})}
 
Вся общая информация инкапсулирована в трёх структурах. Структура
\texttt{RawIMU} хранит одно необработанное измерение:
 
\begin{lstlisting}[caption={Структуры RawIMU и NavState}]
struct RawIMU {
    double  timestamp;
    int16_t xacc,  yacc,  zacc;   // акселерометр (LSB)
    int16_t xgyro, ygyro, zgyro;  // гироскоп (LSB)
    int16_t xmag,  ymag,  zmag;   // магнитометр (LSB)
};
 
struct NavState {
    double qw=1, qx=0, qy=0, qz=0; // кватернион ориентации
    double wx=0, wy=0, wz=0;        // угловые скорости (рад/с)
    double vx=0, vy=0, vz=0;        // скорость (м/с)
    double px=0, py=0, pz=0;        // позиция (м)
    double roll=0, pitch=0, yaw=0;  // углы Эйлера (градусы)
    double ax=0,  ay=0,  az=0;      // ускорение (м/с^2)
};
\end{lstlisting}
 
Структура \texttt{SharedBuffer} объединяет всё межпоточное взаимодействие.
Каждое поле защищено отдельным мьютексом, что позволяет потокам работать
параллельно без глобальной блокировки. Флаг \texttt{running} объявлен как
\texttt{std::atomic<bool>}:
 
\begin{lstlisting}[caption={Потокобезопасный разделяемый буфер}]
struct SharedBuffer {
    std::mutex           mtx;
    std::deque<RawIMU>   queue;         // очередь пакетов (макс. 500)
    static const int     MAX_QUEUE = 500;
    std::atomic<bool>    running{true};
    std::mutex           nav_mtx;
    NavState             nav;
    std::mutex           trail_mtx;
    std::deque<std::array<float,3>> trail; // траектория (макс. 2000)
    static const int     MAX_TRAIL = 2000;
 
    void push(const RawIMU& r) {
        std::lock_guard<std::mutex> lk(mtx);
        if ((int)queue.size() < MAX_QUEUE) queue.push_back(r);
    }
    bool pop(RawIMU& out) {
        std::lock_guard<std::mutex> lk(mtx);
        if (queue.empty()) return false;
        out = queue.front(); queue.pop_front(); return true;
    }
};
\end{lstlisting}
 
\newpage
 
% ─────────────────────────────────────────────────────────────────────────────
\section{Зависимости и система сборки}
% ─────────────────────────────────────────────────────────────────────────────
 
\subsection{Используемые библиотеки}
 
\begin{table}[H]
\centering
\caption{Внешние зависимости проекта}
\label{tab:deps}
\begin{tabularx}{\textwidth}{|l|l|X|}
\hline
\textbf{Библиотека} & \textbf{Версия} & \textbf{Роль в проекте} \\
\hline
MAVLink \texttt{c\_library\_v2}~\cite{mavlink_docs} & 2.0
    & Парсинг телеметрии, декодирование \texttt{RAW\_IMU} \\
\hline
GLFW3~\cite{glfw_docs} & $\geq$3.3
    & Создание окна OpenGL, ввод мыши и клавиатуры \\
\hline
GLEW & $\geq$2.1
    & Загрузка расширений OpenGL~3.3 \\
\hline
GLM  & $\geq$0.9.9
    & Матричная и кватернионная математика \\
\hline
pthreads & POSIX
    & Межпоточная синхронизация (\texttt{std::thread}, \texttt{std::mutex}) \\
\hline
\end{tabularx}
\end{table}
 
Библиотека MAVLink является заголовочной~--- она полностью состоит из
\texttt{.h}-файлов и не требует компиляции~\cite{mavlink_proto}. Остальные
зависимости устанавливаются из репозитория пакетов Ubuntu.
 
\subsection{Конфигурация CMake}
 
Сборочная система использует CMake версии не ниже~3.10 и настроена на
стандарт C++14. Зависимости ищутся через \texttt{find\_package} и
\texttt{PkgConfig}; проект компилируется с флагами \texttt{-O2 -Wall -Wextra}.
CMake ищет GLFW тремя способами последовательно: через \texttt{find\_package},
через \texttt{PkgConfig} и, наконец, вручную через \texttt{find\_path} /
\texttt{find\_library}~--- это обеспечивает совместимость с различными
дистрибутивами Ubuntu.
 
\newpage
 
% ─────────────────────────────────────────────────────────────────────────────
\section{Математика инерциальной навигации}
% ─────────────────────────────────────────────────────────────────────────────
 
\subsection{Калибровка датчиков}
 
Сырые показания МЭМС-датчиков содержат постоянное смещение нуля (bias),
обусловленное технологическими особенностями изготовления
кристалла~\cite{titterton2004}. Смещения определялись экспериментально по
данным CSV-файла в состоянии покоя. Перевод в физические единицы выполняется
линейным масштабированием:
\begin{equation}
    a_i = \bigl(\tilde{a}_i - b_{a,i}\bigr) \cdot S_a, \quad
    \omega_i = \bigl(\tilde{\omega}_i - b_{\omega,i}\bigr) \cdot S_\omega,
    \label{eq:calib}
\end{equation}
где $\tilde{a}_i$, $\tilde{\omega}_i$~--- сырые целочисленные отсчёты;
$b_{a,i}$, $b_{\omega,i}$~--- смещения нуля;
$S_a = 9{,}81/3050$ [м/с$^2$/отсчёт],
$S_\omega = 0{,}001$ [рад/с/отсчёт].
 
Экспериментально определённые смещения, вынесенные в пространство имён
\texttt{Calib} в файле \texttt{imu\_shared.h}:
\begin{equation}
    \mathbf{b}_a = (31,\; 24,\; 0)^\top, \quad
    \mathbf{b}_\omega = (-7,\; 1{,}5,\; 3)^\top.
    \label{eq:biases}
\end{equation}
 
\subsection{Представление ориентации кватернионами}
 
Ориентация тела в пространстве описывается единичным
кватернионом~\cite{kuipers1999}:
\begin{equation}
    \mathbf{q} = q_w + q_x\,\mathbf{i} + q_y\,\mathbf{j} + q_z\,\mathbf{k},
    \quad \|\mathbf{q}\| = 1.
    \label{eq:quat}
\end{equation}
 
Кватернионы предпочтительнее матриц вращения и углов Эйлера: они не страдают
от шарнирного замка (\textit{gimbal lock}), компактны (4~параметра против~9)
и численно стабильны при нормализации~\cite{kuipers1999}.
 
\subsection{Кватернионное интегрирование гироскопа}
 
Кинематическое уравнение связывает производную кватерниона с угловой скоростью
$\bm{\omega} = (\omega_x,\,\omega_y,\,\omega_z)^\top$ в системе координат
тела~\cite{mahony2008}:
\begin{equation}
    \dot{\mathbf{q}} = \tfrac{1}{2}\,\mathbf{q} \otimes
    \begin{pmatrix} 0 \\ \bm{\omega} \end{pmatrix}.
    \label{eq:qdot}
\end{equation}
 
Для численного интегрирования на шаге $\Delta t$ применяется точная формула
через ось-угол~\cite{titterton2004}:
\begin{equation}
    \theta = \|\bm{\omega}\| \cdot \Delta t, \qquad
    \Delta\mathbf{q} = \Bigl(\cos\tfrac{\theta}{2},\;
    \tfrac{\bm{\omega}}{\|\bm{\omega}\|}\sin\tfrac{\theta}{2}\Bigr).
    \label{eq:dq}
\end{equation}
 
Обновление кватерниона с нормализацией:
\begin{equation}
    \mathbf{q} \leftarrow \frac{\mathbf{q} \otimes \Delta\mathbf{q}}
    {\|\mathbf{q} \otimes \Delta\mathbf{q}\|}.
    \label{eq:qupdate}
\end{equation}
 
\subsection{Матрица поворота из кватерниона}
 
Для перевода вектора ускорения из системы координат тела в мировую используется
матрица вращения $\mathbf{R} \in SO(3)$~\cite{kuipers1999}:
\begin{equation}
    \mathbf{R} =
    \begin{pmatrix}
        1{-}2(q_y^2{+}q_z^2) & 2(q_xq_y{-}q_wq_z) & 2(q_xq_z{+}q_wq_y) \\
        2(q_xq_y{+}q_wq_z) & 1{-}2(q_x^2{+}q_z^2) & 2(q_yq_z{-}q_wq_x) \\
        2(q_xq_z{-}q_wq_y) & 2(q_yq_z{+}q_wq_x) & 1{-}2(q_x^2{+}q_y^2)
    \end{pmatrix}.
    \label{eq:rot}
\end{equation}
 
\subsection{Углы Эйлера (конвенция ZYX)}
 
Из кватерниона вычисляются углы крена $\phi$, тангажа $\vartheta$ и рыскания
$\psi$~\cite{titterton2004}:
\begin{align}
    \phi      &= \operatorname{atan2}\!\bigl(2(q_wq_x+q_yq_z),\;
                  1-2(q_x^2+q_y^2)\bigr), \label{eq:roll}\\
    \vartheta &= \arcsin\!\bigl(2(q_wq_y-q_zq_x)\bigr), \label{eq:pitch}\\
    \psi      &= \operatorname{atan2}\!\bigl(2(q_wq_z+q_xq_y),\;
                  1-2(q_y^2+q_z^2)\bigr). \label{eq:yaw}
\end{align}
 
\subsection{Двойное интегрирование (скорость и позиция)}
 
Ускорение переводится из системы тела в мировую, из него вычитается
гравитация~\cite{groves2013}:
\begin{equation}
    \mathbf{a}_\mathrm{world} = \mathbf{R}\,\mathbf{a}_\mathrm{body}
    - \mathbf{g}, \quad
    \mathbf{g} = (0,\;0,\;{-}9{,}81)^\top \;\text{[м/с}^2\text{]}.
    \label{eq:aworld}
\end{equation}
 
Интегрирование методом Эйлера первого порядка:
\begin{align}
    \mathbf{v}(t{+}\Delta t) &= \mathbf{v}(t) +
        \mathbf{a}_\mathrm{world}(t)\cdot\Delta t, \label{eq:vel}\\
    \mathbf{p}(t{+}\Delta t) &= \mathbf{p}(t) +
        \mathbf{v}(t)\cdot\Delta t. \label{eq:pos}
\end{align}
 
\subsection{Метод ZUPT}
 
ZUPT (\textit{Zero Velocity Update})~--- классический способ частичной
коррекции INS в моменты покоя~\cite{skog2010}. Если норма угловой скорости
опускается ниже порогового значения $\omega_\mathrm{th}=0{,}015$~рад/с,
скорость мягко гасится с коэффициентом $\alpha=0{,}85$:
\begin{equation}
    \text{если}\ \|\bm{\omega}\| < \omega_\mathrm{th}:\quad
    \mathbf{v} \leftarrow \mathbf{v}\cdot\alpha.
    \label{eq:zupt}
\end{equation}
Дополнительно, если компонента мирового ускорения по модулю меньше
$0{,}3$~м/с$^2$, она обнуляется полностью, что устраняет интегрирование
шумового сигнала акселерометра в статике.
 
\newpage
 
% ─────────────────────────────────────────────────────────────────────────────
\section{Реализация навигационного алгоритма}
% ─────────────────────────────────────────────────────────────────────────────
 
\subsection{Функция \texttt{integrate\_imu()}}
 
Весь навигационный алгоритм сосредоточен в функции \texttt{integrate\_imu()}
файла \texttt{imu\_visualizer.cpp}. Она вызывается из потока визуализации для
каждого пакета, извлечённого из очереди. Шаг интегрирования $\Delta t$
вычисляется как разность временны'х меток двух последовательных пакетов.
 
\textbf{Блок 1.} Калибровка сырых отсчётов по формуле~\eqref{eq:calib}:
 
\begin{lstlisting}[caption={Калибровка сырых показаний}]
double ax = (raw.xacc  - Calib::AX_BIAS) * Calib::ACCEL_SCALE;
double ay = (raw.yacc  - Calib::AY_BIAS) * Calib::ACCEL_SCALE;
double az = (raw.zacc  - Calib::AZ_BIAS) * Calib::ACCEL_SCALE;
double gx = (raw.xgyro - Calib::GX_BIAS) * Calib::GYRO_SCALE;
double gy = (raw.ygyro - Calib::GY_BIAS) * Calib::GYRO_SCALE;
double gz = (raw.zgyro - Calib::GZ_BIAS) * Calib::GYRO_SCALE;
\end{lstlisting}
 
\textbf{Блок 2.} Кватернионное интегрирование через ось-угол по
формулам~\eqref{eq:dq}--\eqref{eq:qupdate}:
 
\begin{lstlisting}[caption={Кватернионное интегрирование гироскопа}]
double theta = std::sqrt(gx*gx + gy*gy + gz*gz) * dt;
if (theta > 1e-10) {
    double norm_w = std::sqrt(gx*gx + gy*gy + gz*gz);
    double s   = std::sin(theta * 0.5) / norm_w;
    double c   = std::cos(theta * 0.5);
    double dqw = c,    dqx = gx*s, dqy = gy*s, dqz = gz*s;
 
    // q_new = q (x) dq
    double qw = n.qw*dqw - n.qx*dqx - n.qy*dqy - n.qz*dqz;
    double qx = n.qw*dqx + n.qx*dqw + n.qy*dqz - n.qz*dqy;
    double qy = n.qw*dqy - n.qx*dqz + n.qy*dqw + n.qz*dqx;
    double qz = n.qw*dqz + n.qx*dqy - n.qy*dqx + n.qz*dqw;
 
    double norm = std::sqrt(qw*qw + qx*qx + qy*qy + qz*qz);
    n.qw=qw/norm; n.qx=qx/norm; n.qy=qy/norm; n.qz=qz/norm;
}
\end{lstlisting}
 
\textbf{Блок 3.} Перевод ускорения в мировую систему координат,
вычитание гравитации и применение ZUPT:
 
\begin{lstlisting}[caption={Мировое ускорение, компенсация гравитации и ZUPT}]
// Матрица поворота (строки r00..r22) из кватерниона
double ax_w = r00*ax + r01*ay + r02*az;
double ay_w = r10*ax + r11*ay + r12*az;
double az_w = r20*ax + r21*ay + r22*az;
az_w -= (-9.81);         // вычитаем гравитацию
 
// ZUPT: гироскоп в покое -> мягко гасим скорость
double gyro_norm = std::sqrt(gx*gx + gy*gy + gz*gz);
if (gyro_norm < Calib::ZUPT_GYRO_THRESH) {
    n.vx *= 0.85; n.vy *= 0.85; n.vz *= 0.85;
    if (std::abs(ax_w) < 0.3) ax_w = 0;
    if (std::abs(ay_w) < 0.3) ay_w = 0;
    if (std::abs(az_w) < 0.3) az_w = 0;
}
// Двойное интегрирование
n.vx += ax_w*dt; n.vy += ay_w*dt; n.vz += az_w*dt;
n.px += n.vx*dt; n.py += n.vy*dt; n.pz += n.vz*dt;
\end{lstlisting}
 
\subsection{Вычисление углов Эйлера}
 
Завершает функцию вычисление углов Эйлера по
формулам~\eqref{eq:roll}--\eqref{eq:yaw}. Аргумент \texttt{asin}
зажимается в интервал $[-1,1]$ для защиты от накопленной численной
погрешности:
 
\begin{lstlisting}[caption={Углы Эйлера из кватерниона (ZYX)}]
const double R2D = 180.0 / M_PI;
n.roll  = std::atan2(2*(qw*qx + qy*qz),
                     1 - 2*(qx*qx + qy*qy)) * R2D;
double sp = std::max(-1.0, std::min(1.0, 2*(qw*qy - qz*qx)));
n.pitch = std::asin(sp) * R2D;
n.yaw   = std::atan2(2*(qw*qz + qx*qy),
                     1 - 2*(qy*qy + qz*qz)) * R2D;
\end{lstlisting}
 
\newpage
 
% ─────────────────────────────────────────────────────────────────────────────
\section{Реализация визуализации}
% ─────────────────────────────────────────────────────────────────────────────
 
\subsection{Шейдеры OpenGL~3.3}
 
Визуализатор использует две пары GLSL-шейдеров. Первая
(\texttt{VERT3D}/\texttt{FRAG3D}) применяется для рендеринга 3D-модели
полётного контроллера с диффузным освещением по модели Ламберта:
 
\begin{lstlisting}[language={},caption={Вершинный и фрагментный шейдеры 3D-модели}]
// VERT3D
#version 330 core
layout(location=0) in vec3 aPos;
layout(location=1) in vec3 aColor;
layout(location=2) in vec3 aNormal;
out vec3 vColor; out vec3 vNormal; out vec3 vPos;
uniform mat4 uMVP; uniform mat4 uModel;
void main(){
    vPos    = vec3(uModel * vec4(aPos, 1));
    vNormal = mat3(transpose(inverse(uModel))) * aNormal;
    vColor  = aColor;
    gl_Position = uMVP * vec4(aPos, 1);
}
 
// FRAG3D
#version 330 core
in vec3 vColor; in vec3 vNormal; in vec3 vPos;
out vec4 outColor;
uniform vec3 uLight; uniform float uAlpha;
void main(){
    float d = max(dot(normalize(vNormal), normalize(uLight)), 0.0);
    outColor = vec4(vColor * (0.35 + 0.65*d), uAlpha);
}
\end{lstlisting}
 
Вторая пара (\texttt{VERT\_FLAT}/\texttt{FRAG\_FLAT}) используется для
отрисовки всех линейных примитивов без освещения: координатной сетки,
мировых осей, траектории и HUD-графиков.
 
\subsection{Процедурная модель полётного контроллера}
 
3D-модель FC строится процедурно из параллелепипедов функцией
\texttt{build\_fc()}. Вспомогательная \texttt{push\_box()} генерирует
шесть граней куба с различными коэффициентами яркости для имитации объёма;
\texttt{push\_arrow()} добавляет три цветных стрелки осей (X~--- красная,
Y~--- зелёная, Z~--- синяя):
 
\begin{lstlisting}[caption={Процедурная сборка модели FC}]
static std::vector<V3> build_fc() {
    std::vector<V3> v;
    // Корпус -- тёмно-серый прямоугольник
    push_box(v,  0,     0,    0,    0.8f,0.8f,0.14f, 0.13f,0.14f,0.19f);
    // Белая полоска "нос" (направление Y+)
    push_box(v,  0,  0.37f, 0.08f, 0.7f,0.06f,0.04f, 0.95f,0.95f,0.95f);
    // Боковые рёбра
    push_box(v,  0.37f, 0,  0.08f, 0.04f,0.8f,0.04f, 0.22f,0.22f,0.28f);
    push_box(v, -0.37f, 0,  0.08f, 0.04f,0.8f,0.04f, 0.22f,0.22f,0.28f);
    // Стрелки осей
    push_arrow(v, 0,0,0.1f,  1,0,0, 0.7f, 0.95f,0.10f,0.10f); // X
    push_arrow(v, 0,0,0.1f,  0,1,0, 0.7f, 0.10f,0.92f,0.10f); // Y
    push_arrow(v, 0,0,0.07f, 0,0,1, 0.7f, 0.15f,0.42f,0.98f); // Z
    return v;
}
\end{lstlisting}
 
\subsection{Главный цикл рендера}
 
Функция \texttt{run\_visualizer()} инициализирует контекст GLFW/GLEW,
компилирует шейдеры, загружает геометрию в видеопамять и запускает главный
цикл. Схема каждого кадра:
\begin{enumerate}[leftmargin=2cm]
    \item Обработать события ввода (мышь, клавиши).
    \item Вычитать из \texttt{SharedBuffer} \textit{все} накопившиеся
          пакеты, для каждого вызвать \texttt{integrate\_imu()}, добавить
          точку траектории и обновить кольцевые буферы.
    \item Обновить матрицы орбитальной камеры, следящей за FC.
    \item Отрисовать \textbf{3D-сцену} (верхние 72\% экрана):
          сетку, мировые оси, траекторию с цветовым затуханием, модель FC.
    \item Отрисовать \textbf{HUD-панель} (нижние 30\% экрана):
          четыре графика временных рядов.
\end{enumerate}
 
Матрица модели FC формируется из текущей позиции и кватерниона ориентации с
помощью GLM:
 
\begin{lstlisting}[caption={Матрица модели из позиции и кватерниона}]
glm::mat4 fc_pos = glm::translate(glm::mat4(1),
    glm::vec3((float)nav.px, (float)nav.py, (float)nav.pz));
glm::quat  q((float)nav.qw, (float)nav.qx,
             (float)nav.qy, (float)nav.qz);
glm::mat4 fc_model = fc_pos * glm::mat4_cast(q);
glm::mat4 fc_mvp   = vp * fc_model;
\end{lstlisting}
 
Траектория отрисовывается как набор отрезков с цветовым градиентом: старые
точки темнее, новые~--- ярче:
 
\begin{lstlisting}[caption={Траектория с цветовым затуханием}]
float fade_step = 1.0f / buf.trail.size();
for (int i = 0; i < (int)buf.trail.size()-1; i++) {
    float alpha = (float)i * fade_step;
    float cr = 0.1f + 0.7f * alpha;  // красный: 0.1 -> 0.8
    float cg = 0.5f + 0.4f * alpha;  // зелёный: 0.5 -> 0.9
    float cb = 0.9f;                  // синий: константа
    // ... добавить отрезок trail[i] -> trail[i+1]
}
\end{lstlisting}
 
\subsection{HUD-панель графиков}
 
Нижняя панель содержит четыре независимых графика на основе скользящего
буфера \texttt{RingBuf} (300~последних значений). Каждый график отображает
три сигнала (по осям X, Y, Z) красным, зелёным и синим цветами:
 
\begin{table}[H]
\centering
\caption{Содержимое HUD-панелей}
\label{tab:hud}
\begin{tabularx}{\textwidth}{|>{\centering\arraybackslash}p{3cm}|
                               >{\centering\arraybackslash}p{3cm}|
                               X|}
\hline
\textbf{Панель} & \textbf{Диапазон} & \textbf{Сигналы} \\
\hline
Акселерометр & $\pm15$~м/с$^2$ & $a_x$ (кр.), $a_y$ (зел.), $a_z$ (синий) \\
\hline
Гироскоп     & $\pm0{,}15$~рад/с & $\omega_x$, $\omega_y$, $\omega_z$ \\
\hline
Углы Эйлера  & $\pm180\degree$   & крен, тангаж, рыскание \\
\hline
Скорость     & $\pm2$~м/с        & $v_x$, $v_y$, $v_z$ \\
\hline
\end{tabularx}
\end{table}
 
\newpage
 
% ─────────────────────────────────────────────────────────────────────────────
\section{Инструкция по установке и запуску}
% ─────────────────────────────────────────────────────────────────────────────
 
\subsection{Установка зависимостей}
 
\begin{lstlisting}[language=bash,caption={Установка системных зависимостей (Ubuntu)}]
sudo apt update
sudo apt install build-essential cmake git \
    libglfw3-dev libglew-dev libglm-dev libgl1-mesa-dev
\end{lstlisting}
 
Библиотека MAVLink является заголовочной и клонируется напрямую в каталог
проекта:
 
\begin{lstlisting}[language=bash,caption={Клонирование MAVLink}]
cd pixhawk_imu_visualizer
git clone https://github.com/mavlink/c_library_v2.git
\end{lstlisting}
 
\subsection{Сборка проекта}
 
\begin{lstlisting}[language=bash,caption={Сборка через CMake}]
mkdir build && cd build
cmake ..
make -j$(nproc)
\end{lstlisting}
 
При успешной сборке в каталоге \texttt{build/} появляется исполняемый файл
\texttt{pixhawk\_imu\_reader}.
 
\subsection{Настройка прав и запуск}
 
\begin{lstlisting}[language=bash,caption={Права на порт и запуск}]
# Добавить пользователя в группу dialout (необходима перезагрузка)
sudo usermod -a -G dialout $USER
# Или временно:
sudo chmod 666 /dev/ttyACM0
 
cd build
./pixhawk_imu_reader
\end{lstlisting}
 
Окно визуализации открывается немедленно. Для завершения~--- Ctrl+C в
терминале.
 
\subsection{Управление в окне визуализации}
 
\begin{table}[H]
\centering
\caption{Управление в окне визуализации}
\label{tab:controls}
\begin{tabularx}{\textwidth}{|l|X|}
\hline
\textbf{Действие} & \textbf{Эффект} \\
\hline
ЛКМ + перетаскивание & Вращение камеры вокруг объекта \\
\hline
Колесо мыши          & Приближение / отдаление (0{,}5~---~30~м) \\
\hline
Клавиша \textbf{R}   & Сброс позиции и траектории в $(0,\,0,\,0)$ \\
\hline
\textbf{Ctrl+C}      & Завершение программы \\
\hline
\end{tabularx}
\end{table}
 
\newpage
 
% ─────────────────────────────────────────────────────────────────────────────
\section{Результаты работы}
% ─────────────────────────────────────────────────────────────────────────────
 
\subsection{Внешний вид программы}
 
На рисунке~\ref{fig:main} представлен общий вид окна визуализатора в
реальном времени. Верхняя область содержит 3D-сцену с
моделью полётного контроллера на фоне координатной сетки и мировых осей;
нижняя --- четыре графика датчиков.
 
\begin{figure}[H]
    \centering

    \includegraphics[width=\linewidth]{figures/screen_main}
    
    \caption{Главное окно визуализатора: 3D-сцена и HUD-панель}
    \label{fig:main}
\end{figure}
 
На рисунке~\ref{fig:rotate} показано состояние визуализатора при вращении
полётного контроллера: 3D-модель следует за реальной ориентацией устройства,
а кривые гироскопа на HUD-панели фиксируют угловые скорости по каждой из осей.
 
\begin{figure}[H]
    \centering
   
    \includegraphics[width=\linewidth]{figures/screen_rotate}
    
    \caption{Визуализатор при вращении устройства: реакция HUD-графиков гироскопа}
    \label{fig:rotate}
\end{figure}
 
На рисунке~\ref{fig:hud} приведён укрупнённый вид HUD-панели. Нулевые линии
в каждом графике выделены серым. Видно, что в статике (покой платформы) кривые
гироскопа сходятся к нулю~--- это подтверждает корректность подобранных
смещений~\eqref{eq:biases}.
 
\begin{figure}[H]
    \centering
    
    \includegraphics[width=\linewidth]{figures/screen_hud}
    
    \caption{Укрупнённый вид HUD-панели: акселерометр, гироскоп, углы Эйлера, скорость}
    \label{fig:hud}
\end{figure}
 
\subsection{Результаты оценки ориентации}
 
В ходе тестирования подтверждена корректная работа алгоритма кватернионного
интегрирования. При вращении полётного контроллера по всем трём осям
3D-модель точно воспроизводила движения устройства. Углы Эйлера, вычисленные
по формулам~\eqref{eq:roll}--\eqref{eq:yaw}, соответствовали физическим
поворотам.
 
\subsection{Результаты оценки позиции}
 
Алгоритм двойного интегрирования ускорения показал неудовлетворительные
результаты при оценке абсолютных координат. Ключевая проблема~---
\textbf{дрейф позиции}: уже через 5--10~с накопленная ошибка достигала
нескольких метров. Механизм ZUPT снизил скорость нарастания дрейфа в
состоянии покоя, однако не устранил его полностью из-за ненулевого шума
акселерометра порядка $10^{-3}$--$10^{-2}$~м/с$^2$.
 
\subsection{Анализ причин дрейфа}
 
Пусть истинное ускорение $a(t)$ и измеренное $\hat{a}(t)$ связаны:
\begin{equation}
    \hat{a}(t) = a(t) + b + \eta(t),
    \label{eq:noise}
\end{equation}
где $b$~--- остаточное смещение после калибровки, $\eta(t)$~--- случайный
шум. Ошибка позиции после двойного интегрирования за время $T$:
\begin{equation}
    \delta p(T) = \tfrac{1}{2}\,b\,T^2 +
    \int_0^T\!(T-t)\,\eta(t)\,dt.
    \label{eq:drift}
\end{equation}
 
Первое слагаемое растёт \textit{квадратично} по времени даже при малом
остаточном смещении~\cite{groves2013}. За 10~с это даёт ошибку порядка
нескольких метров~\cite{woodman2007}.
 
\subsection{Возможные направления развития}
 
Полноценная инерциальная навигация без внешних корректирующих источников
принципиально невозможна~\cite{groves2013}. Среди возможных улучшений:
\begin{enumerate}[leftmargin=2cm]
    \item \textbf{Фильтр Калмана с GPS-коррекцией}~\cite{groves2013}~---
          классический подход для авиационных INS.
    \item \textbf{Фильтр Маджвика}~\cite{madgwick2011}~--- градиентный
          алгоритм слияния данных гироскопа, акселерометра и магнитометра.
    \item \textbf{Нелинейный дополняющий фильтр на SO(3)}~\cite{mahony2008}.
    \item \textbf{ZUPT с фильтром Калмана}~\cite{skog2010}~---
          расширенное применение метода нулевой скорости.
\end{enumerate}
 
\newpage
 
% ─────────────────────────────────────────────────────────────────────────────
\section*{Заключение}
\addcontentsline{toc}{section}{Заключение}
% ─────────────────────────────────────────────────────────────────────────────
 
В ходе выполнения курсовой работы разработана система визуализации
инерциальной навигации на основе полётного контроллера Pixhawk.
Реализованы следующие компоненты:
\begin{itemize}[leftmargin=2cm]
    \item автономный модуль первичного считывания \texttt{RAW\_IMU} с записью
          данных в CSV для анализа и подбора калибровочных коэффициентов;
    \item многопоточная архитектура C++14 с потокобезопасным
          \texttt{SharedBuffer} на основе \texttt{std::mutex} и
          \texttt{std::atomic};
    \item алгоритм калибровки сырых показаний с экспериментально
          определёнными смещениями нуля~\eqref{eq:biases};
    \item точное кватернионное интегрирование угловой скорости через
          ось-угол~\cite{kuipers1999} с нормализацией на каждом шаге;
    \item вычисление матрицы вращения $\mathbf{R}\in SO(3)$ и углов Эйлера
          из кватерниона;
    \item двойное интегрирование ускорения для оценки позиции~\cite{groves2013};
    \item механизм ZUPT~\cite{skog2010} для частичной компенсации дрейфа;
    \item OpenGL~3.3-визуализация: процедурная 3D-модель FC, координатная
          сетка, мировые оси, траектория с цветовым затуханием, HUD-панель
          с четырьмя графиками~\cite{glfw_docs}.
\end{itemize}
 
Алгоритм оценки ориентации показал высокое качество: 3D-модель точно
воспроизводила реальные повороты устройства. Алгоритм оценки позиции
продемонстрировал принципиальное ограничение МЭМС-датчиков~\cite{titterton2004}:
ошибка нарастает квадратично по времени~\eqref{eq:drift}, и получение
практически пригодных координат требует слияния данных IMU с корректирующими
сенсорами в рамках фильтра Калмана~\cite{groves2013}.
 
\newpage
 
% ─────────────────────────────────────────────────────────────────────────────
\section*{Список использованных источников}
\addcontentsline{toc}{section}{Список использованных источников}
% ─────────────────────────────────────────────────────────────────────────────
 
\begin{thebibliography}{99}
 
\bibitem{titterton2004}
    Titterton~D., Weston~J.
    \textit{Strapdown Inertial Navigation Technology}.
    2nd~ed.~--- IET, 2004.~--- 558~p.
 
\bibitem{woodman2007}
    Woodman~O.\,J.
    \textit{An Introduction to Inertial Navigation}.
    Technical Report UCAM-CL-TR-696.~---
    University of Cambridge, 2007.~--- 37~p.~---
    URL: \url{https://www.cl.cam.ac.uk/techreports/UCAM-CL-TR-696.pdf}
    (дата обр.~01.03.2024).
 
\bibitem{kuipers1999}
    Kuipers~J.\,B.
    \textit{Quaternions and Rotation Sequences}.~---
    Princeton University Press, 1999.~--- 398~p.
 
\bibitem{mahony2008}
    Mahony~R., Hamel~T., Pflimlin~J.-M.
    Nonlinear Complementary Filters on the Special Orthogonal Group~//
    \textit{IEEE Transactions on Automatic Control}.~---
    2008.~--- Vol.~53, No.~5.~--- P.~1203--1218.~---
    DOI:~10.1109/TAC.2008.923738.
 
\bibitem{madgwick2011}
    Madgwick~S.\,O.\,H., Harrison~A.\,J.\,L., Vaidyanathan~R.
    Estimation of IMU and MARG orientation using a gradient descent
    algorithm~//
    \textit{IEEE International Conference on Rehabilitation Robotics
    (ICORR)}.~--- 2011.~--- P.~1--7.~---
    DOI:~10.1109/ICORR.2011.5975346.
 
\bibitem{groves2013}
    Groves~P.\,D.
    \textit{Principles of GNSS, Inertial, and Multisensor Integrated
    Navigation Systems}.
    2nd~ed.~--- Artech House, 2013.~--- 776~p.
 
\bibitem{skog2010}
    Skog~I., Handel~P.
    In-Car Positioning and Navigation Technologies~--- A Survey~//
    \textit{IEEE Transactions on Intelligent Transportation Systems}.~---
    2009.~--- Vol.~10, No.~1.~--- P.~4--21.~---
    DOI:~10.1109/TITS.2008.2011712.
 
\bibitem{mavlink_proto}
    Meier~L. et al.
    MAVLink: A Lightweight Message Marshalling Library for the
    DroneCode Ecosystem.~---
    URL: \url{https://github.com/mavlink/mavlink}
    (дата обр.~01.03.2024).
 
\bibitem{mavlink_docs}
    MAVLink Developer Guide: Common Messages.~---
    URL: \url{https://mavlink.io/en/messages/common.html}
    (дата обр.~01.03.2024).
 
\bibitem{px4_docs}
    PX4 Autopilot User Guide.~---
    URL: \url{https://docs.px4.io/main/en/}
    (дата обр.~05.03.2024).
 
\bibitem{glfw_docs}
    GLFW~--- An OpenGL library. Documentation.~---
    URL: \url{https://www.glfw.org/docs/latest/}
    (дата обр.~05.03.2024).
 
\bibitem{posix_termios}
    The Open Group. The Single UNIX Specification, Version~4.
    \textit{termios.h --- define values for termios}.~---
    URL: \url{https://pubs.opengroup.org/onlinepubs/9699919799/basedefs/termios.h.html}
    (дата обр.~01.03.2024).
 
\end{thebibliography}
 
\end{document}
